\begin{proof}[Proof of \cref{obj:lem:measurable-kernel-radon-nikodym}]
\begin{enumerate}
\item Use the canonical kernel Radon--Nikodym derivative to define
\[
f:H\times Z\longrightarrow[0,\infty],
\qquad
f(h,z):=\frac{dK}{dL}(h,z).
\]
The kernel Radon--Nikodym construction gives joint measurability of this function for finite kernels with standard-Borel target. % lean: ProbabilityTheory.Kernel.measurable_rnDeriv

\item Fix \(h\in H\) and a measurable set \(E\subseteq Z\). The canonical kernel derivative agrees with the ordinary measure derivative:
\[
f(h,z)
=
\frac{dK(h,\cdot)}{dL(h,\cdot)}(z)
\qquad\text{for \(L(h,\cdot)\)-almost every \(z\)}.
\]
Consequently,
\[
\int_E f(h,z)\,L(h,dz)
=
\int_E
\frac{dK(h,\cdot)}{dL(h,\cdot)}(z)\,L(h,dz)
=
K(h,E).
\]
The first equality uses the almost-everywhere agreement, and the second is the measure Radon--Nikodym identity, applicable because \(K(h,\cdot)\ll L(h,\cdot)\) and both measures are finite. Reversing this equality gives the required identity for every \(h\) and \(E\). % lean: ProbabilityTheory.Kernel.setLIntegral_rnDeriv
\end{enumerate}
\end{proof}
